% TOP_PROBLEM: {"id": 30006807, "title": "Witten's Asymptotic Expansion Conjecture for SU(2) WRT Invariants", "category_id": 7, "category": {"id": 7, "name": "topology", "display_name": "Topology", "description": "Properties preserved under continuous deformations.", "slug": "topology", "order_index": 7, "created_at": "2026-07-31T15:26:25.671Z"}, "status": "open"}
% FIRST_POSTED: 2026-10-01
% ATTEMPT_STATUS: unresolved
% !TeX program = lualatex
% Definition and sources only; this file is not an LLM solution attempt.
\documentclass[11pt]{article}
\usepackage[margin=25mm]{geometry}
\usepackage{fontspec}
\setmainfont{Latin Modern Roman}
\usepackage{amsmath,amssymb,mathtools}
\usepackage{unicode-math}
\setmathfont{Latin Modern Math}
\usepackage{xurl,hyperref}
\hypersetup{colorlinks=true,urlcolor=blue}
\setcounter{secnumdepth}{0}
\setlength{\parindent}{0pt}
\setlength{\parskip}{5pt}
\setlength{\emergencystretch}{3em}
\begin{document}
\section*{\texorpdfstring{Witten's Asymptotic Expansion Conjecture for SU(2) WRT Invariants}{Witten's Asymptotic Expansion Conjecture for SU(2) WRT Invariants}}\label{30006807}
\subsection{Definitions and mathematical statement}
For a closed oriented three-manifold $Y$, write $\mathrm{WRT}_k(Y)$ for the normalized $\mathrm{SU}(2)$ quantum invariant at level $k-2$, with value one on $S^3$. Flat $\mathrm{SU}(2)$ connections have Chern--Simons actions $S\in{\mathbb{R}}/{\mathbb{Z}}$; denote their set of values by $\mathrm{CS}(Y)$. The target is a full Poincar\'e expansion
\[
 \mathrm{WRT}_k(Y)\sim\sum_{S\in\mathrm{CS}(Y)}e^{2\pi i kS}W_S(1/k),
 \qquad W_S(\tau)\in\bigcup_{n\ge1}{\mathbb{C}}((\tau^{1/n})).
\]
A Puiseux Laurent series has rational-power exponents with one finite denominator and only finitely many negative powers. Full asymptotics requires a remainder estimate after every prescribed truncation, not only the leading order. The separate conjectures specifying growth exponents or leading coefficients are not added here.
\par\smallskip\textit{Formulation and concept references:} \href{https://arxiv.org/html/2510.10678v1}{[S1]}.

\subsection{Short English statement}
Do SU(2) quantum invariants of every closed three-manifold have complete large-level expansions assembled from its flat-connection Chern–Simons phases?
\subsection{Research attempt: product stability of phased asymptotic expansions}
\textbf{Status.} We prove an analytic closure lemma relevant to connected sums and verify the normalized sphere case. Applying the lemma to manifolds requires the corresponding exact topological-quantum-field-theory product identity; the proof below states that input explicitly. It does not prove the asymptotic expansion for arbitrary prime three-manifolds. The primary source [S1] treats a Seifert family, rather than all closed three-manifolds.

\subsubsection{1. A precise multiplication lemma}
Suppose two sequences admit full expansions
\[A_k\sim\sum_{s\in S}e^{2\pi iks}a_s(k^{-1}),\qquad
B_k\sim\sum_{t\in T}e^{2\pi ikt}b_t(k^{-1}),\]
where $S,T$ are finite and each coefficient series is Puiseux Laurent with a finite denominator. Then their product has the expansion
\[A_kB_k\sim\sum_{u\in S+T}e^{2\pi iku}
\sum_{s+t=u}a_s(k^{-1})b_t(k^{-1}),\tag{1}\]
where addition of phases is modulo one. A common denominator is the least common multiple of the finitely many denominators. The sum of two lower bounds on Laurent exponents bounds the product exponents below, so every coefficient in each product series is a finite sum. Coincident phases must be combined, as shown, and cancellation may make a coefficient series zero.

Here is the remainder argument, which formal multiplication alone would omit. Full expansions imply polynomial growth $A_k=O(k^a)$ and $B_k=O(k^b)$ for some nonnegative $a,b$. For any requested error $O(k^{-R})$, choose finite truncations $A_k^{\le},B_k^{\le}$ with errors respectively $O(k^{-R-b})$ and $O(k^{-R-a})$, increasing both truncation depths if necessary. The truncations have the same polynomial growth bounds. Then
\[A_kB_k-A_k^{\le}B_k^{\le}
=(A_k-A_k^{\le})B_k+A_k^{\le}(B_k-B_k^{\le})=O(k^{-R}).\]
Deleting terms of the finite product below the requested accuracy changes the result by another $O(k^{-R})$. Since $R$ is arbitrary, (1) is a full Poincar\'e expansion. Truncating both factors at the same naive order without accounting for growth could fail this estimate.

\subsubsection{2. Conditional connected-sum propagation}
In a three-dimensional TQFT whose sphere state space is one-dimensional, let the invariant of $S^3$ be a nonzero scalar $z$. Removing a ball from a closed manifold produces a state in that one-dimensional space. The gluing pairing therefore implies the unnormalized formula $Z(M\#N)=Z(M)Z(N)/z$: evaluate both states as multiples of the ball state, whose self-pairing is $z$. Dividing all closed-manifold invariants by $z$ gives an invariant $I$ with $I(S^3)=1$ and
\[I(M\#N)=I(M)I(N).\tag{2}\]
Thus, whenever the target WRT normalization is realized by this standard gluing convention, (2) and the multiplication lemma propagate the target expansion from $M,N$ to their connected sum.

The phase condition has a geometric counterpart. The fundamental group of a connected sum is the free product of the summand groups. A representation into $\mathrm{SU}(2)$ is specified by its two restrictions. Flat connections can be trivialized near the balls and glued across a product neck; the Chern--Simons integral adds, with changes of trivialization affecting only the integer ambiguity. Consequently the resulting phases add modulo one. Under this standard gluing description, the phases in (1) are permitted target phases. This conclusion uses actual flat-connection gluing, not arbitrary multiplication of phase lists.

\subsubsection{3. Exact basic case and the scope of the reduction}
By the selected normalization $\mathrm{WRT}_k(S^3)=1$. Its fundamental group is trivial, so the only flat-connection phase is zero, and the expansion has $W_0=1$ with zero remainder at every order. The lemma also propagates any independently established family of expansions through the connected-sum operation under (2). This includes no claim that an unknown summand can be canceled from a product expansion: a factor can have zeros or small values, and asymptotic division requires additional hypotheses.

\subsubsection{4. Audit and remaining gap}
Finite phase sets and finite Puiseux denominators are used at the exact places where common denominators and finite coefficient sums are formed. If logarithms or other terms were needed for an unknown manifold, this closure argument would not rule them out. It proves no stationary-phase theorem for singular character varieties and supplies no all-orders bounds for arbitrary prime manifolds. Those analytic tasks remain the missing core of the full conjecture.

\subsection{Sources}
\begin{itemize}
\item[S1]Jorgen Ellegaard Andersen, Li Han, Yong Li, William Elbaek Mistegard, David Sauzin and Shanzhong Sun, A proof of Witten's asymptotic expansion conjecture for WRT invariants of Seifert fibered homology spheres (2025).
\url{https://arxiv.org/html/2510.10678v1}.
\textit{Formulation source.}
\end{itemize}
\end{document}
